%! Linear Combinations of Vectors — Unit 1, Lesson 1.1 — Student Packet Cover Sheet
\documentclass[10pt]{article}
\usepackage{linalg-article}
\usepackage{linalg-boxes}
\usepackage{ltablex}
\keepXColumns

\begin{document}

% Full-bleed burgundy banner
\begin{tikzpicture}[remember picture, overlay]
  \fill[burgundy] (current page.north west)
    rectangle ([yshift=-0.9in]current page.north east);
\end{tikzpicture}
\vspace{-0.8in}

\begin{center}
  {\color{white}\textbf{\LARGE Linear Algebra}}\\[4pt]
  {\color{blushmid}\large\textbf{Unit 1 \quad Vectors and Matrices}}\\[6pt]
  {\color{white}\textbf{\large Lesson 1.1 \quad Linear Combinations of Vectors}}
\end{center}

\vspace{0.22in}
\namedateperiod
\vspace{0.18in}

\begin{learningtargetbox}
\small
\begin{enumerate}[leftmargin=1.5em, itemsep=5pt]
  \item \ldots{} build a \textbf{linear combination} $c\,\mathbf{v}+d\,\mathbf{w}$ by scaling each vector and then adding.
  \item \ldots{} describe what the \textbf{multiples of one vector} reach (a line) and what \textbf{two vectors} together reach (the plane).
  \item \ldots{} find the \textbf{weights} $c,d$ that combine two vectors into a given target, and explain when a target is out of reach.
\end{enumerate}
\end{learningtargetbox}

\vspace{0.14in}

\begin{tocbox}
\small
\renewcommand{\arraystretch}{1.5}
\begin{tabularx}{\linewidth}{@{} c l X r @{}}
\rowcolor{blush}
\textbf{\#} & \textbf{Component} & \textbf{Description} & \textbf{Score} \\
\midrule
1 & Warm-Up        & Spiral review: scaling, adding, and multiples of a vector   & \blank{1.2cm} \\
2 & Guided Notes   & Combinations $c\mathbf{v}+d\mathbf{w}$; lines, planes, and weights & \blank{1.2cm} \\
3 & Group Activity & Robot Moves --- reaching targets by combining vectors       & \blank{1.2cm} \\
4 & Exit Ticket    & Quick check on combinations and span                        & \blank{1.2cm} \\
5 & Homework       & Practice with combinations, weights, and a recipe extension & \blank{1.2cm} \\
\midrule
  & \multicolumn{2}{r}{\textbf{Total}} & \blank{1.2cm} \\
\end{tabularx}
\end{tocbox}

\vspace{0.10in}

\begin{remindbox}
\small A \textbf{linear combination} of $\mathbf{v}$ and $\mathbf{w}$ is
$c\,\mathbf{v}+d\,\mathbf{w}$ --- scale each by a number, then add. The multiples of a single
vector fill a \textbf{line} through the origin; the combinations of two \emph{non-parallel}
vectors fill the whole \textbf{plane}.
\end{remindbox}

\end{document}
